% ════════════════════════════════════════════════════════════════════
% Chapter 6 — Host-Side Software Architecture
% ════════════════════════════════════════════════════════════════════
\section{Host-Side Software Architecture}\label{sec:hostsw}

The host-side application is a C++17 real-time data acquisition, visualisation, and session management platform.
It is built with CMake~3.20+ and depends on: \code{librealsense2}, OpenCV~4.0+, Eigen3, \code{nlohmann/json}, \code{spdlog}, Dear~ImGui~1.90.4 with ImPlot~0.16, GLFW~3.3, and OpenGL.
The application version at the time of writing is v1.1.0.


\subsection{Software Architecture Overview}

The application follows a layered architecture with clear separation of concerns:

\begin{table}[H]
  \centering
  \caption{Host application software layers and key classes.}
  \label{tab:sw_layers}
  \begin{tabular}{llp{6cm}}
    \toprule
    \textbf{Layer} & \textbf{Key Classes} & \textbf{Responsibility} \\
    \midrule
    \textit{Sensors}    & \class{IMUReader}, \class{CameraReader}, \class{MarkerTracker}
                         & Hardware abstraction; thread-safe streaming. \\
    \textit{Core}       & \class{Session}, \class{SyncEngine}, \class{DataLogger}, \class{CalibrationManager}, \class{EventLog}
                         & Session lifecycle, clock mapping, data persistence, calibration. \\
    \textit{Processing} & \class{OrientationFilter}, \class{RepSegmenter}, \class{Validator}, \class{Autoregulation}
                         & Real-time kinematics, rep counting, plausibility checks. \\
    \textit{GUI}        & \class{MainWindow}, \class{SessionPanel}, \class{PreflightPanel}, \class{OperatorView}, etc.
                         & ImGui-based multi-panel tiled interface. \\
    \textit{Annotation}  & \class{AnnotationStudio}, \class{TimelinePanel}, \class{VideoPanel}, \class{RepTablePanel}
                         & Post-session NLE-style timeline editing and validation. \\
    \textit{Utils}      & \class{SerialPort}, \class{AudioCue}, \class{Notifications}, \class{RingBuffer}, \class{SubjectRegistry}
                         & Cross-platform I/O, audio cues, toast system. \\
    \bottomrule
  \end{tabular}
\end{table}


\subsection{Sensor Layer}

\subsubsection{IMUReader}
The \class{IMUReader} class manages a dedicated background thread that continuously reads from the USB serial port at 921\,600~baud.
Incoming bytes are appended to a 4\,KB ring buffer.
The parser scans for the sync word (\code{0x55AA}), extracts 26-byte packets, validates the CRC16-CCITT checksum, and converts raw integers to physical units using the scale factors in \cref{eq:accel_scale}--\cref{eq:temp_scale}.

Quality metrics are continuously accumulated:
\begin{itemize}[nosep]
  \item \textbf{CRC error count:} Packets with failed checksums are discarded and counted.
  \item \textbf{Dropout detection:} Gaps in the ESP32 timestamp sequence exceeding \SI{2}{\milli\second} are counted as dropouts.
  \item \textbf{Jitter statistics:} A rolling window of inter-sample timestamp deltas is maintained to compute mean, max, and standard deviation of sampling jitter.
  \item \textbf{Saturation detection:} Samples at the $\pm$FSR boundaries ($\pm 32768$ raw) are counted per axis, flagging potential clipping events.
  \item \textbf{Noise floor estimation:} Rolling standard deviation of gyroscope magnitude and $|a| - 1$ at rest, providing a live noise-floor proxy.
  \item \textbf{FSYNC hit tracking:} The \code{fsync\_tagged} flag is set when \code{(temp\_raw \& 0x01) == 1}, and a cumulative FSYNC hit counter and live rate are maintained.
\end{itemize}

Each validated sample is passed to a user-registered callback, which the \class{Session} class uses to feed the \class{SyncEngine}, \class{DataLogger}, \class{OrientationFilter}, and \class{RepSegmenter}.

\subsubsection{CameraReader}
The \class{CameraReader} wraps the \code{librealsense2} pipeline.
It configures the D455 according to \cref{tab:cam_config} and starts a background thread that calls \code{pipeline.wait\_for\_frames()} in a blocking loop.
Each frameset is decomposed into:
\begin{itemize}[nosep]
  \item Left and right IR images (\code{cv::Mat}, \code{CV\_8U}).
  \item Aligned depth image (\code{cv::Mat}, \code{CV\_16U}, millimetres).
  \item Hardware timestamp (via \code{RS2\_FRAME\_METADATA\_SENSOR\_TIMESTAMP}).
  \item Frame number (monotonic counter).
\end{itemize}

Camera frames are enqueued into a bounded producer--consumer queue (\code{CAMERA\_QUEUE\_LIMIT = 512}).
A separate camera worker thread dequeues frames and performs marker tracking, video encoding, and data logging, ensuring that the RealSense streaming thread is never blocked by downstream processing.

When enabled, the D455's onboard BMI085 IMU streams accelerometer (\SI{250}{\hertz}) and gyroscope (\SI{200}{\hertz}) data via a separate callback.
These samples are logged to \file{camera\_imu.csv} for tripod-shake detection and cross-stream synchronisation validation, but are never used for VBT kinematic computation.


\subsection{Session Lifecycle}\label{sec:session}

The \class{Session} class orchestrates the entire recording workflow through a well-defined state machine:

\begin{figure}[H]
  \centering
  \begin{tikzpicture}[
    state/.style={draw, rectangle, rounded corners=5pt, minimum width=2.5cm, minimum height=1cm, align=center, font=\small},
    >={Stealth[length=2.5mm]},
    ]
    \node[state] (idle)       {IDLE};
    \node[state, right=2cm of idle] (conf) {CONFIGURED};
    \node[state, right=2cm of conf] (calib) {CALIBRATING};
    \node[state, below=1.5cm of calib] (ready) {READY};
    \node[state, left=2cm of ready] (rec) {RECORDING};
    \node[state, left=2cm of rec] (stop) {STOPPED};
    \node[state, below=1.5cm of stop] (saved) {SAVED};

    \draw[->] (idle) -- node[above, font=\scriptsize] {create()} (conf);
    \draw[->] (conf) -- node[above, font=\scriptsize] {calibrate} (calib);
    \draw[->] (calib) -- node[right, font=\scriptsize] {tap-test OK} (ready);
    \draw[->] (ready) -- node[above, font=\scriptsize] {start\_recording()} (rec);
    \draw[->] (rec) -- node[above, font=\scriptsize] {stop\_recording()} (stop);
    \draw[->] (stop) -- node[left, font=\scriptsize] {save()} (saved);
    \draw[->, dashed] (stop.south) -- ++(0,-0.5) -| node[near start, below, font=\scriptsize] {discard()} (idle.south);
  \end{tikzpicture}
  \caption{Session state machine.}
  \label{fig:session_states}
\end{figure}

\subsubsection{Atomic Session Finalisation}
During recording, all files are written to a directory suffixed with \code{.partial/}.
Upon calling \code{save()}, the directory is atomically renamed to its final path.
If a final directory already exists (e.g., from a previous aborted save), it is moved to a timestamped \code{.bak\_*} backup.
This guarantees that a crash or power failure mid-recording never corrupts the dataset---the \code{.partial} directory remains as a recoverable artefact.

On startup, \code{Session::find\_orphaned\_partials()} scans the dataset root for any unfinished \code{.partial} directories and surfaces them via toast notification, offering the operator the option to recover or discard.

\subsubsection{Multi-Set Support}
Schema~v4 introduces multi-set sessions.
During a continuous recording, the operator can call \code{advance\_set(next\_set\_info)}, which:
\begin{enumerate}[nosep]
  \item Stamps \code{t\_end\_unified\_s} on the current \code{SetInfo} and records the completed rep count.
  \item Appends a new \code{SetInfo} to the session's \code{sets} vector.
  \item Instructs the \class{RepSegmenter} to tag all subsequent reps with the new \code{set\_id}.
\end{enumerate}


\subsection{Data Logging}

The \class{DataLogger} writes data in parallel to multiple output streams:

\begin{table}[H]
  \centering
  \caption{Output files generated per session.}
  \label{tab:output_files}
  \begin{tabular}{llp{5.5cm}}
    \toprule
    \textbf{Directory} & \textbf{File} & \textbf{Content} \\
    \midrule
    \file{imu/}     & \file{raw\_imu.csv}          & All columns: \code{esp\_timestamp\_us}, \code{host\_timestamp\_s}, \code{unified\_time\_s}, \code{accel\_\{x,y,z\}\_g}, \code{gyro\_\{x,y,z\}\_dps}, \code{temp\_c}, raw integers, \code{fsync\_tagged}. \\
    \file{imu/}     & \file{raw\_imu.bin}           & 26-byte binary packets (identical to firmware format). \\
    \file{imu/}     & \file{camera\_imu.csv}        & D455 onboard BMI085 accel/gyro data. \\
    \file{camera/}  & \file{marker\_positions.csv}  & Per-frame: timestamp, 2D pixel coords, 3D world coords, blob area, circularity, SNR, confidence, depth source. \\
    \file{camera/}  & \file{ir\_video.mp4}          & Left IR frames encoded as H.264 video. \\
    \file{camera/}  & \file{video\_frames.csv}      & Per-frame index: HW timestamp, host timestamp, frame number. \\
    \file{annotations/} & \file{rep\_segments.json}  & Per-rep phase boundaries and kinematic summaries. \\
    Root            & \file{session\_info.json}     & Complete session metadata (see \cref{sec:dataset}). \\
    Root            & \file{manifest.json}          & Per-file SHA-256 checksums and sizes. \\
    Root            & \file{events.jsonl}           & Append-only event log with hash chain. \\
    \bottomrule
  \end{tabular}
\end{table}


\subsection{Real-Time Processing}

\subsubsection{Orientation Filter}
The real-time GUI uses a Madgwick AHRS filter (\class{OrientationFilter}) for live gravity-compensated acceleration preview and coordinate-frame visualisation.
The filter gain ($\beta = 0.1$) is tuned for responsiveness during live operation.
Post-processing in the golden model pipeline replaces this with VQF for superior static accuracy.

\subsubsection{Repetition Segmenter}
The \class{RepSegmenter} implements a dual-algorithm approach:
\begin{itemize}[nosep]
  \item \textbf{Algorithm A (Primary):} Camera-derived velocity zero-crossing detection. A 2nd-order Butterworth lowpass filter (cutoff from the exercise profile, e.g., \SI{10}{\hertz} for squats) is applied to the camera velocity signal. A windowed peak-confirmation detector (\code{PEAK\_WINDOW\_N = 20} samples, $\pm$\SI{0.22}{\second} at \SI{90}{fps}) identifies confirmed TOP and BOTTOM extrema with prominence thresholds, avoiding false triggers from noise and oscillation.
  \item \textbf{Algorithm B (Fallback):} Raw IMU accelerometer variance detection. When camera tracking is lost for $>$\SI{1}{\second}, the segmenter switches to monitoring the running acceleration variance over a \SI{200}{\milli\second} window. Motion onset and cessation are detected via variance thresholds.
\end{itemize}

Each completed repetition is stored as a \code{RepAnnotation} containing four chronological phase segments: \textbf{concentric} $\rightarrow$ \textbf{top\_rest} $\rightarrow$ \textbf{eccentric} $\rightarrow$ \textbf{rest}, with associated peak velocity, mean velocity, and ROM.


\subsection{Calibration System}

The \class{CalibrationManager} supports:
\begin{enumerate}[nosep]
  \item \textbf{6-Position Accelerometer Calibration:} The operator places the IMU in six orientations ($\pm X$, $\pm Y$, $\pm Z$ up). For each position, the mean acceleration vector is recorded over $\sim$\SI{3}{\second}. A least-squares ellipsoid fit computes the $3\times3$ scale-misalignment matrix $\mathbf{K}$ and the $3\times1$ bias vector $\mathbf{b}$:
  \begin{equation}
    \mathbf{a}_\text{corrected} = \mathbf{K} \left( \mathbf{a}_\text{raw} - \mathbf{b} \right)
  \end{equation}
  \item \textbf{Per-Session Gyroscope Bias:} The first \SI{5}{\second} of stationary data are averaged to compute the gyroscope bias vector, which is subtracted from all subsequent samples.
  \item \textbf{Camera-to-World Extrinsic:} Gravity alignment from the D455's onboard accelerometer computes the camera-to-world rotation matrix.
\end{enumerate}

All calibration parameters are serialised to JSON with timestamps and sensor serial numbers, forming the \code{CalibrationProvenance} block in the session metadata.


\subsection{Build Provenance}

Every build embeds a \file{Version.h} generated by CMake from \file{Version.h.in}.
This captures:
\begin{itemize}[nosep]
  \item Application version (currently v1.1.0).
  \item Git SHA, branch name, and dirty flag.
  \item Build timestamp (ISO~8601 UTC).
  \item Build type (Release/Debug), compiler identification, OS, and architecture.
\end{itemize}

This fingerprint is embedded into every session's \code{BuildProvenance} metadata block, ensuring that any kinematic result can be traced back to the exact software state that produced it.


\subsection{Graphical User Interface}

The GUI is implemented using Dear~ImGui with ImPlot for real-time charting.
Key panels include:

\begin{itemize}[nosep]
  \item \textbf{Pre-Flight Panel:} A systematic checklist that gates the START button. Checks include: IMU rate within $\pm 5\%$ of \SI{1}{\kilo\hertz}, no accelerometer saturation, camera connected, marker detection rate $> 80\%$, sync drift $<$ \SI{50}{ppm}, and subject ID populated.
  \item \textbf{Operator View (F12):} A fullscreen, big-numbers display readable from across the gym. Shows: current rep count, last MCV and PCV, velocity-loss percentage, and a stop-set recommendation when VL exceeds the exercise profile threshold.
  \item \textbf{Rep Timeline Panel:} A horizontal NLE-style timeline showing per-rep bars with colour-coded plausibility dots (green = passed, red = failed). Supports drag-handle boundary editing, right-click deletion, and bulk operations.
  \item \textbf{Calibration Wizard (F1):} A step-by-step 6-position guide with live alignment visualisation (dot product + magnitude check), 3-second countdown capture, and a review screen.
  \item \textbf{Replay Mode (F2):} Browse completed sessions, load metadata, view rep annotations, inspect the manifest and event log.
\end{itemize}
